% ECONOMICS_PROBLEM: {"id": "C81-341", "title": "Measurement error in digital and administrative data", "jel_code": "C81", "source_id": "OP-0108"}
% !TeX program = xelatex
\documentclass[11pt,a4paper]{article}
\usepackage[margin=23mm]{geometry}
\usepackage{fontspec}
\setmainfont{Latin Modern Roman}
\usepackage{amsmath,amssymb,mathtools}
\usepackage{xcolor,enumitem,booktabs,longtable,array,xurl,hyperref}
\definecolor{muted}{HTML}{53636C}
\hypersetup{colorlinks=true,urlcolor=blue,linkcolor=blue}
\setlength{\parindent}{0pt}
\setlength{\parskip}{5pt}
\newcommand{\R}{\mathbb R}
\newcommand{\E}{\mathbb E}
\newcommand{\N}{\mathbb N}
\newcommand{\ind}{\mathbf 1}
\newcommand{\Var}{\operatorname{Var}}
\newcommand{\Cov}{\operatorname{Cov}}
\newcommand{\supp}{\operatorname{supp}}
\DeclareMathOperator*{\argmin}{arg\,min}
\DeclareMathOperator*{\argmax}{arg\,max}
\newcommand{\entrymeta}[3]{{\sffamily\small\color{muted}\textbf{\texttt{#1}}\quad|\quad #2\par #3\par}}
\newcommand{\Problemref}[1]{\texttt{#1}}
\newcommand{\JELref}[2]{\texttt{#2} (JEL \texttt{#1})}
\begin{document}
\section*{Measurement error in digital and administrative data}
\textbf{Problem ID:} \texttt{C81-341}\quad \textbf{JEL:} \texttt{C81}\par
% BEGIN ECONOMICS STATEMENT
\label{problem:OP-0108}
\label{atlas:prob:E8}

\noindent\textbf{Original identifier:} E8 (atlas item 108). \textbf{Field:} Econometrics and causal inference. \textbf{Classification:} editorial research family with established restricted solutions; current open status not certified.

\subsection*{Definitions and assumptions}

Let latent research variables $L=(Y,D,X)$ follow a population law, with $D\in\{0,1\}$, bounded outcome $Y\in[a,b]$, and covariates $X$. Operational data report $W=(\widetilde Y,\widetilde D,\widetilde X)$ according to a measurement kernel $K_s(dw\mid L)$ indexed by system version $s$. A validation indicator $V$ reveals $L$ when $V=1$, in addition to $W,s$. Specify validation sampling probabilities $v(W,s)$, and assume validation is independent of $L$ conditional on $(W,s)$ only if defensible. Under latent-variable consistency, unconfoundedness and overlap, the causal target is $\theta=\mathbb E[m_1(X)-m_0(X)]$, $m_d(x)=\mathbb E[Y\mid D=d,X=x]$.

\subsection*{Formal research question}

Characterize the identified set for $\theta$ and estimators uniformly robust to nonclassical measurement and system changes. In a representative validation design with $v(W,s)\ge\eta>0$, the joint latent-observed law is recoverable in principle; quantify the information needed when the validation fraction tends to zero. With no validation for a new version $s$, instead impose a declared kernel distance $d(K_s,K_{s_0})\le\delta$ from a validated version $s_0$, where $d$ and $\delta$ are supplied. Find sharp causal-effect bounds over all latent laws and kernels matching the observed law and these restrictions, and confidence sets covering the bounds uniformly as sample size, validation size and drift vary.

\subsection*{Interpretation and scope}

An algorithmic score is not inherently a noisy measurement of the desired construct: the construct and intervention must first be defined. Record linkage can create mismatched outcomes, selective missingness and dependent observations, so a linkage error rate alone may not characterize bias. Multiple proxies identify latent variables only with additional conditional-independence or structural restrictions; more observations cannot identify an unrestricted measurement kernel. Classical errors-in-variables attenuation holds in particular linear models and should not be assumed for misclassified treatment or outcomes. Correcting measurement does not remove unobserved causal confounding. Provenance identifies the versioning structure and validation design on which credible comparisons depend.

\subsection*{Source audit and status}

[S1] and [S2] are verified measurement-error references, covering classical and nonclassical concerns. The third input item, \emph{modern record-linkage literature}, is a field label and cannot be audited as a publication; it is retained as unresolved [S3]. The formal kernel-and-validation question is an editorial specialization. The references motivate it but do not supply a universal modern-digital-data theorem or a certified 2026 open-status claim.

\subsection*{References}

\begin{enumerate}[label={[\textnormal{S}\arabic*]},leftmargin=*]

\item Jerry Hausman (2001). \emph{Mismeasured Variables in Econometric Analysis: Problems from the Right and Problems from the Left}. Journal of Economic Perspectives 15(4), 57–67. \url{https://doi.org/10.1257/jep.15.4.57}. \textbf{Locator:} Publisher or institutional abstract and bibliographic record. \textbf{Support and limits:} Measurement-error overview; classical attenuation needs specific assumptions.

\item John Bound, Charles Brown and Nancy Mathiowetz (2001). \emph{Measurement Error in Survey Data}. Handbook of Econometrics 5, chapter 59, 3705–3843. \url{https://www.sciencedirect.com/science/article/pii/S1573441201050127}. \textbf{Locator:} Publisher or institutional abstract and bibliographic record. \textbf{Support and limits:} Classical and nonclassical measurement, validation and proxies; not platform-specific digital measurement.

\item \textbf{Unresolved original citation:} modern record-linkage literature. Not an identifiable publication: no authors, title, year or persistent identifier supplied.

\end{enumerate}

\subsection*{Common conventions: Source mathematical conventions}

Unless an entry states otherwise, agent and firm sets are finite. State and action spaces are measurable, strategies and outcome maps are measurable, probability laws are specified, and expectations exist for the targets under discussion. Infinite-horizon discount factors lie in $(0,1)$ and discounted objectives are finite. Norms, tolerances and complexity input sizes must be specified for computational comparisons. Constants or primitives that appear locally are inputs, not empirical facts asserted by this catalogue.

\textbf{Identification.} A statistical model $m\in\mathcal M$ implies an observable law $P_m$ and target $\theta(m)$. At law $P$, its identified set is
\[
\Theta(P;\mathcal M)=\{\theta(m):m\in\mathcal M,\ P_m=P\}.
\]
Point identification holds when this set is a singleton, or equivalently when $P_{m_1}=P_{m_2}$ implies $\theta(m_1)=\theta(m_2)$ for all compatible models. Sharp bounds mean bounds attained, or approached when the boundary is not attained, by compatible models. Writing this set is a definition, not a solution: the substantive task is to establish informative restrictions and derive or estimate the set. An unrestricted-confounding model may give only trivial bounds.

\textbf{Inference.} Identification, estimability and valid uncertainty quantification are separate questions. Fix $\alpha\in(0,1)$. A confidence procedure $C_n$ over a declared law class $\mathcal P$ has asymptotic uniform coverage at level $1-\alpha$ when
\[
\liminf_{n\to\infty}\inf_{P\in\mathcal P}\Pr_P\{\theta(P)\in C_n\}\geq1-\alpha.
\]
For set-identified targets the coverage event must be defined separately, for example coverage of the full identified set. If an entry uses identification-indexed models instead of uniquely indexed laws, the same coverage convention is applied over those models. Pointwise limits do not establish uniform validity, and asymptotic validity does not guarantee finite-sample precision. Sample sizes, dependence structures and weak-identification sequences are part of each application.

\textbf{Counterfactuals.} $Y_i(a)$ is a potential outcome under a specified intervention, including its timing, implementation and versions. It is not $Y_i$ conditional on voluntarily choosing $a$. A full intervention vector permits interference; shorthand scalar treatment notation does not silently impose no interference. Assignment, selection, attrition, anticipation and measurement must be specified. A claim about an instrument's local effect is not automatically a population policy effect.

\textbf{Economic equilibrium.} A counterfactual model includes best responses, constraints, feasibility, market clearing, state transitions and expectations consistent with the stated information. When several equilibria exist, either a declared selector is an input or the target is set-valued. Comparisons across policy regimes must use the stated selection convention. Reduced-form relationships are not presumed invariant after an equilibrium-changing intervention.

\textbf{Optimization and welfare.} Optimization is over the stated feasible set. If attainment is not established, $\sup$ or $\inf$ and $\varepsilon$-optimal choices replace an asserted maximizer or minimizer. For example, with value $V=\sup_{a\in\mathcal A}W(a)<\infty$, an $\varepsilon$-optimal set is $\{a\in\mathcal A:W(a)\geq V-\varepsilon\}$ for $\varepsilon>0$. Benefits, costs, risk and health or environmental outcomes can be aggregated only using declared units and conversion weights. Interpersonal utility weights, the treatment of fiscal transfers and foreign profits, discounting, and the behavioral welfare criterion are explicit inputs. A comparative-static derivative additionally requires existence and appropriate differentiability of the selected counterfactual.

\textbf{Sources.} Local labels [S1], [S2], and so forth restart in every question. Locators identify the relevant abstract, model, section or result at the level actually checked; a bibliographic or abstract-level verification is not represented as inspection of an entire proof. Working papers are distinguished from later publications and changing versions. Source summaries are paraphrased. All URL checks and editorial formulations refer to the audit date stated above.
% END ECONOMICS STATEMENT
\end{document}
